%%%%%%%%%%%%%%%%%%%%%%%%%%%%%%%%%%%%%%%%%%%%%%%%%%%%%%%%%%%%%%%%%%%%%%
% Lớp 11 | Hàm số lượng giác và phương trình lượng giác | Các công thức lượng giác
% Dạng 1D1?3-3: Áp dụng công thức nhân đôi - hạ bậc
% Số câu: 45  (sinh tự động bởi tach_cau.py; từ nay thầy sửa trực tiếp file này)
%%%%%%%%%%%%%%%%%%%%%%%%%%%%%%%%%%%%%%%%%%%%%%%%%%%%%%%%%%%%%%%%%%%%%%

%%%%%%%%%%%%%%%% TRẮC NGHIỆM 4 PHƯƠNG ÁN %%%%%%%%%%%%%%%%

% nguồn: Ok_Toan11_TracNghiem_Lan2
\begin{ex}%[1D1N3-3]%[ma:83a8c4fc]
	Với $\alpha$ là số thực bất kỳ, trong các mệnh đề sau,  mệnh đề nào \textbf{sai}?
	\choice 
		{$\sin 2\alpha =2\sin \alpha \cdot \cos  \alpha$}
		 {$\cos 2\alpha =2\cos^2\alpha-1$}
		 {$\cos 2\alpha =-2\sin^2\alpha+1$}
		{\True $\cos 2\alpha=\sin^2\alpha-\cos^2\alpha$} 
	 \loigiai{
		Ta có $\sin 2\alpha=2\sin \alpha  \cdot \cos \alpha $; $\cos 2\alpha =\cos^2\alpha -\sin^2\alpha =2\cos^2\alpha-1=1-2\sin^2\alpha$.
	 }
\end{ex}

% nguồn: Ok_Toan11_TracNghiem_Lan2
\begin{ex}%[1D1H3-3]%[ma:8d083f77]
	Cho $\sin 2\alpha =\dfrac {3}{4}$. Tính giá trị biểu thức  $A=\tan \alpha +\cot \alpha$.
	\choice 
		{$A=\dfrac  {4}{3}$}
		{$A=\dfrac {2}{3}$} 
		 {\True $A=\dfrac {8}{3}$}
		 {$A=\dfrac {16}{3}$}
	\loigiai{ 
		\allowdisplaybreaks 
		\begin{align*}
			A&=\tan \alpha +\cot \alpha\\ 
			&=\dfrac {\sin \alpha}{\cos \alpha}+\dfrac {\cos \alpha}{\sin \alpha}\\ 
			&=\dfrac {\sin^2\alpha +\cos^2\alpha }{\sin \alpha \cos  \alpha}\\
			&=\dfrac  {1}{\dfrac {1}{2}\sin 2\alpha}\\
			&=\dfrac {1}{\dfrac {1}{2}\cdot \dfrac {3}{4}}\\ 
			&=\dfrac  {8}{3}.
		\end{align*}
	 }
\end{ex}

% nguồn: Ok_Toan11_TracNghiem_Lan2
\begin{ex}%[1D1H3-3]%[ma:763d871e]
	Khẳng định nào  dưới đây \textbf{sai}?
	\choice 
		{$2\sin  ^2a=1-\cos 2a$}
		{\True $\cos 2a=2\cos  a-1$}
		{$\sin 2a=2\sin a\cos a$} 
		{$\sin \left(a+b\right)=\sin a\cos b+\sin  b\cos a$}
	 \loigiai{
		 Có $\cos 2a=2\cos ^2a-1$.
	} 
\end{ex}

% nguồn: Ok_Toan11_TracNghiem_Lan2
\begin{ex}%[1D1H3-3]%[ma:c2a96eab]
	Cho $a$, $b$ là hai góc nhọn. Biết $\cos  a=\dfrac {1}{3}$, $\cos b=\dfrac {1}{4}$. Giá trị của biểu thức \\$\cos (a+b) \cos (a-b)$ bằng
	\choice 
		{\True $-\dfrac {119}{144}$} 
		 {$-\dfrac {115}{144}$}
		{$-\dfrac {113}{144}$} 
		{$-\dfrac {117}{144}$} 
	\loigiai{Từ $\cos a=\dfrac  {1}{3}\Rightarrow \cos 2a=2\cos^2a-1=-\dfrac {7}{9}$.\\
		$\cos b=\dfrac  {1}{4}\Rightarrow \cos 2b=2\cos^2b-1=-\dfrac {7}{8}$.\\
		Ta có $\cos \left(a+b\right)\cos \left(a-b\right)=\dfrac {1}{2}\left( \cos 2a+\cos 2b \right)=\dfrac {1}{2}\left(-\dfrac  {7}{9}-\dfrac {7}{8} \right)=-\dfrac {119}{144}$.
	 }
\end{ex}

% nguồn: Ok_Toan11_TracNghiem_Lan2
\begin{ex}%[1D1H3-3]%[ma:f69c05d4]
	Biểu  thức $\dfrac{3-4\cos 2\alpha+\cos 4\alpha}{3+4\cos 2\alpha+\cos 4\alpha}$ có kết quả rút gọn bằng
	\choice 
		 {$-\tan ^4\alpha$}
		{\True $\tan ^4\alpha$} 
		 {$-\cot ^4\alpha$}
		 {$\cot ^4\alpha$}
	 \loigiai{
		\begin{eqnarray*}
			\dfrac{3-4\cos 2\alpha+\cos 4\alpha}{3+4\cos 2\alpha+\cos 4\alpha}&=&\dfrac{3-4\left(1-2\sin ^2\alpha \right)+2\left(1-2\sin ^2\alpha \right)^2-1}{3+4\left(2\cos ^2\alpha-1\right)+2\left(2\cos ^2\alpha-1\right)^2-1}  \\
			&  =& \dfrac{8\sin ^2a-8\sin ^2\alpha+8\sin ^4\alpha}{8\cos ^2a-8\cos ^2\alpha+8\cos ^4\alpha}\\
			 &=&\tan ^4\alpha.
		\end{eqnarray*}
	} 
\end{ex}

% nguồn: Ok_Toan11_TracNghiem_Lan2
\begin{ex}%[1D1H3-3]%[ma:db9fa1f0]
	Rút gọn biểu thức $A=\dfrac{\sin x+\sin \dfrac{x}{2}}{1+\cos x+\cos \dfrac{x}{2}}$ được kết quả nào sau  đây?
	 \choice
		{\True $\tan \dfrac{x}{2}$} 
		{$\cot  x$}
		 {$\tan ^2\left(\dfrac{\pi}{4}-x\right)$}
		{$\sin x$} 
	\loigiai{ 
		Ta có 
		\begin{eqnarray*}
			 A & = & \dfrac{2 \sin \dfrac{x}{2}\cos \dfrac{x}{2} +\sin \dfrac{x}{2}}{ 2 \cos^2 \dfrac{x}{2}+\cos \dfrac{x}{2}}\\
			& = & \dfrac{\sin \dfrac{x}{2} \cdot \left( 2  \cos\dfrac{x}{2}+ 1\right)}{\cos \dfrac{x}{2}\cdot \left( 1 + 2 \cos \dfrac{x}{2}\right)}\\
			& = & \tan  \dfrac{x}{2}.
		\end{eqnarray*}
	} 
\end{ex}

% nguồn: Ok_Toan11_TracNghiem_Lan2
\begin{ex}%[1D1H3-3]%[ma:6f0e90cb]
	Cho  $\sin 2\alpha=-\dfrac{4}{5}$ và $\dfrac{3\pi}{4}<\alpha<\pi$. Giá trị của $\sin \alpha$ là
	 \choice
		 {$\dfrac {2}{5}$}
		{$\dfrac {1}{5}$} 
		{$\dfrac  {2\sqrt {5}}{5}$}
		{\True $\dfrac {\sqrt {5}}{5}$} 
	 \loigiai{
		 Ta có $\dfrac {3\pi }{4}<\alpha<\pi \Rightarrow \sin \alpha>0$.\\
		$\dfrac {3\pi }{2}<2\alpha<2\pi \Rightarrow  \cos 2\alpha>0$.\\
		$\cos^2 2\alpha =1-\sin^22\alpha =1-\left(-\dfrac {4}{5}\right)^2=\dfrac {9}{25}\Rightarrow \cos 2\alpha =\dfrac  {3}{5}$.\\
		$\sin^2\alpha =\dfrac {1-\cos 2\alpha}{2}=\dfrac {1-\dfrac {3}{5}}{2}=\dfrac {1}{5}\Rightarrow \sin \alpha=\dfrac  {\sqrt {5}}{5}$.
	} 
\end{ex}

% nguồn: Ok_Toan11_TracNghiem_Lan2
\begin{ex}%[1D1H3-3]%[ma:b3da5427]
	Biết rằng $\sin 18^\circ=\dfrac {a+b\sqrt {5}}{c}$, với $a$, $b$, $c\in \mathbb{Z}$, $c\ne 0$ và $\dfrac {a}{c}$, $\dfrac {b}{c}$ là các phân số tối giản. Giá trị của biểu thức $S=a+b+c$  là
	 \choice
		{\True  $S=2$}
		 {$S=4$}
		{$S=3$} 
		{$S=1$} 
	\loigiai{ 
		Ta  có
		 \allowdisplaybreaks
		\begin{eqnarray*}
			& &\cos  36^\circ =\sin 54^\circ\\
			&\Leftrightarrow &  1-2\sin^218^\circ =3\sin 18^\circ -4\sin^318^\circ\\
			&\Leftrightarrow & 4\sin^318^\circ -2\sin^218^\circ  -3\sin 18^\circ +1=0 \\
			&\Leftrightarrow  & (\sin 18^\circ -1)(4\sin ^218^\circ +2\sin 18^\circ -1)=0\\
			&\Leftrightarrow & \hoac{&\sin 18^\circ =1\\&\sin 18^\circ =\dfrac {-1\pm \sqrt  {5}}{2}.}
		\end{eqnarray*}
		Vì $0^\circ <18^\circ <90^\circ$ nên $0<\sin 18^\circ <1$,  do đó $\sin 18^\circ =\dfrac {-1+\sqrt {5}}{2}$.\\
		Suy ra $a=-1$,  $b=1$, $c=2$.\\
		Vậy  $S=a+b+c=2$.
	 }
\end{ex}

% nguồn: Ok_Toan11_TracNghiem_Lan2
\begin{ex}%[1D1H3-3]%[ma:45565bf5]
	Công thức nào sau  đây \textbf{sai}?
	 \choice
		{$\cos 6a = \cos^2  3a - \sin^2 3a$}
		{$\cos 6a = 1  - 2\sin^2 3a$}
		 {\True $\cos 6a = 1 - 6\sin^2 a$}
		{$\cos 6a = 2\cos^2  3a - 1$}
	\loigiai{ 
		Áp dụng công thức $\cos 2\alpha  = \cos^2 \alpha - \sin^2 \alpha = 2\cos^2 \alpha - 1 = 1 - 2\sin^2 \alpha$, ta được\\
		 $\cos 6a = \cos^2 3a - \sin^2 3a= 2\cos^2 3a - 1= 1 - 2\sin^2 3a$.
	} 
\end{ex}

% nguồn: Ok_Toan11_TracNghiem_Lan2
\begin{ex}%[1D1H3-3]%[ma:96c136c9]
	Chọn  đẳng thức đúng.
	\choice 
		{\True $\cos^2\left(\dfrac{\pi}{4} +  \dfrac{a}{2}\right) = \dfrac{1 - \sin a}{2}$}
		{$\cos^2\left(\dfrac{\pi}{4} + \dfrac{a}{2}\right) =  \dfrac{1 + \sin a}{2}$}
		 {$\cos^2\left(\dfrac{\pi}{4} + \dfrac{a}{2}\right) = \dfrac{1 - \cos a}{2}$}
		{$\cos^2\left(\dfrac{\pi}{4} + \dfrac{a}{2}\right) = \dfrac{1 + \cos  a}{2}$}
	\loigiai{Ta có 
		$\cos^2 \left( \dfrac{\pi}{4} + \dfrac{a}{2} \right) = \dfrac{1 + \cos \left( \dfrac{\pi}{2} + a \right)}{2} = \dfrac{1 + \sin(-a)}{2} = \dfrac{1 - \sin  a}{2}$.
	 }
\end{ex}

% nguồn: Ok_Toan11_TracNghiem_Lan2
\begin{ex}%[1D1H3-3]%[ma:aadc236a]
	Cho số thực $\alpha$ thỏa mãn $\sin \alpha =\dfrac  {1}{4}$. Tính $\left(\sin 4\alpha +2\sin 2\alpha \right) \cos \alpha$.
	 \choice
		 {$\dfrac {25}{128}$}
		{$\dfrac {1}{16}$} 
		 {$\dfrac {255}{128}$}
		{\True  $\dfrac {225}{128}$}
	 \loigiai{Ta có
		\allowdisplaybreaks 
		\begin{align*}
			\left(\sin 4\alpha +2\sin 2\alpha \right)\cos \alpha &=2\sin 2\alpha \left(\cos 2\alpha +1 \right)\cos \alpha\\ 
			&=4\sin \alpha  \cos \alpha \left(1-2\sin^2\alpha +1\right)\cos \alpha\\
			&=4\sin \alpha \left( 1-\sin ^2\alpha \right)\left(2-2\sin^2\alpha  \right)\\
			&=8\left(1-\sin^2\alpha \right)^2\sin \alpha\\ 
			&=  8\left(1-\dfrac {1}{16} \right)^2\cdot \dfrac {1}{4}\\
			 &=\dfrac {225}{128}.
		\end{align*}
	} 
\end{ex}

% nguồn: Ok_Toan11_TracNghiem_Lan2
\begin{ex}%[1D1H3-3]%[ma:e08de21c]
	Biết rằng $\sin^6x+\cos^6x=a+b\sin^22x$, với  $a$, $b$ là các số thực. Tính $T=3a+4b$.
	\choice 
		{$T=-7$} 
		{$T=1$} 
		{\True  $T=0$}
		 {$T=7$}
	 \loigiai{Ta có
		 \allowdisplaybreaks
		\begin{align*}
			\sin^6x+\cos^6x&=\left(\sin^2x+\cos^2x\right)^3-3\sin^2x\cdot \cos^2x\left(\sin^2x+\cos^2x\right)\\ 
			&=1-3\sin^2x\cdot  \cos^2x\\
			 &=1-\dfrac {3}{4}\sin^22x.
		\end{align*}
		Vậy $a=1$, $b=-\dfrac {3}{4}$. Do đó $T=3a+4b=0$. 
	} 
\end{ex}

% nguồn: Ok_Toan11_TracNghiem_Lan2
\begin{ex}%[1D1H3-3]%[ma:4406c22d]
	Nếu $\tan \dfrac{x}{2}=\dfrac{1}{2}$ thì giá trị của biểu thức  $\dfrac{\sin x}{2-3\cos x}$ bằng
	 \choice
		{$1$} 
		{$2$} 
		 {$3$}
		{\True  $4$}
	 \loigiai{
		\begin{eqnarray*}
			\dfrac{\sin x }{2 -3 \cos x}&=& \dfrac{2 \sin \dfrac{x}{2}\cos \dfrac{x}{2}}{ 2 - 3 \cos ^2  \dfrac{x}{2}+ 3 \sin^2 \dfrac{x}{2}}\\
			& =& \dfrac{2 \tan  \dfrac{x}{2}}{ 2 \left( 1+\tan^2 \dfrac{x}{2}\right) - 3 +3 \tan^2 \dfrac{x}{2}}\\&=& \dfrac{2 \cdot \dfrac{1}{2}}{ 2 \left(1 + \dfrac{1}{4} \right) - 3 + 3 \cdot \dfrac{1}{4}} = 4.
		\end{eqnarray*}
	 }
\end{ex}

% nguồn: Ok_Toan11_TracNghiem_Lan2
\begin{ex}%[1D1H3-3]%[ma:af67d531]
	Tính giá trị  của biểu thức $M = \cos 10^\circ \cos 20^\circ \cos 40^\circ \cos 80^\circ$.
	 \choice
		{$M=\dfrac{1}{16}\cos  10^\circ$}
		{$M=\dfrac{1}{2}\cos  10^\circ$}
		{$M=\dfrac{1}{4}\cos 10^\circ$} 
		{\True $M=\dfrac{1}{8}\cos 10^\circ$} 
	\loigiai{ 
		Vì $\sin 10^\circ \ne  0$ nên suy ra
		\begin{eqnarray*}
			&M=&\dfrac{16\sin 10^\circ \cos 10^\circ  \cos 20^\circ \cos 40^\circ \cos 80^\circ}{16\sin 10^\circ} \\
			&M=&\dfrac{8\sin 20^\circ  \cos 20^\circ \cos 40^\circ \cos 80^\circ}{16\sin 10^\circ}\\
			&M=&\dfrac{4\sin 40^\circ \cos 40^\circ \cos  80^\circ}{16\sin 10^\circ}\\
			&M=&\dfrac{2\sin 80^\circ \cos 80^\circ}{16\sin 10^\circ}\\ 
			&M=&\dfrac{\sin 160^\circ}{16\sin 10^\circ}\\ 
			&M=&\dfrac{\sin  20^\circ}{16\sin 10^\circ}\\
			&M=&\dfrac{2\sin 10^\circ  \cos 10^\circ}{16\sin 10^\circ}\\
			 &M=&\dfrac{1}{8}\cos 10^\circ.
		\end{eqnarray*}
	 }
\end{ex}

% nguồn: Ok_Toan11_TracNghiem_Lan2
\begin{ex}%[1D1H3-3]%[ma:6f12128f]
	Cho $\cot a=15$, giá trị $\sin  2a$ có thể nhận giá trị nào dưới đây
	\choice 
		{$\dfrac {11}{113}$} 
		{$\dfrac {13}{113}$} 
		 {\True $\dfrac {15}{113}$}
		{$\dfrac {17}{113}$} 
	\loigiai{$\cot a=15 \Rightarrow \dfrac {1}{\sin^2a}=226\Rightarrow \heva{& \sin^2a=\dfrac {1}{226}\\&\cos^2a=\dfrac {225}{226}}\Rightarrow  \sin 2a=\pm \dfrac {15}{113}$.
	 }
\end{ex}

% nguồn: Ok_Toan11_TracNghiem_Lan2
\begin{ex}%[1D1H3-3]%[ma:c800dfad]
	Nếu $\tan  \dfrac{x}{2}=\dfrac{a}{b}$ thì biểu thức $a\sin x+b\cos x$ bằng
	\choice 
		{$a$} 
		{\True $b$} 
		{$\dfrac{a+b}{a}$} 
		{$\dfrac{a+b}{b}$} 
	\loigiai{ 
		Ta có $\tan \dfrac{x}{2}= \dfrac{a}{b}$  nên $ \sin \dfrac{x}{2}= \dfrac{a}{b} \cdot \cos \dfrac{x}{2} $.
		\\ 
		Khi  đó:
		\begin{eqnarray*}
			a \sin  x + b \cos x &=& 2 a \sin \dfrac{x}{2} \cos \dfrac{x}{2} - 2b \sin ^2 \dfrac{x}{2}+ b \\
			&=& 2 \dfrac{a^2}{b} \cos ^2 \dfrac{x}{2} -2 b \cdot \dfrac{a^2}{b^2}\cdot \cos^2 \dfrac{x}{2}+b \\ 
			& =  & b.
		\end{eqnarray*}
	} 
\end{ex}

% nguồn: Ok_Toan11_TracNghiem_Lan2
\begin{ex}%[1D1H3-3]%[ma:286c1975]
	Nếu $\sin  x+\cos x=\dfrac {1}{2}$ thì $\sin2x$ bằng
	 \choice
		{$\dfrac {3}{4}$} 
		{$\dfrac {3}{8}$} 
		 {$\dfrac {\sqrt {2}}{2}$}
		{\True  $-\dfrac {3}{4}$}
	 \loigiai{Ta có $\sin x+\cos x=\dfrac {1}{2}\Leftrightarrow \sin^2x+2\sin x\cos x+\cos^2x=\dfrac {1}{4}$ $\Leftrightarrow \sin 2x=-\dfrac {3}{4}$.
	 }
\end{ex}

% nguồn: Ok_Toan11_TracNghiem_Lan2
\begin{ex}%[1D1H3-3]%[ma:da345ded]
	Rút gọn biểu thức $A  = \dfrac{\tan \alpha - \cot \alpha}{\tan \alpha + \cot \alpha} + \cos 2\alpha$.
	\choice 
		{\True $0$} 
		{$2 \cos^2 x$} 
		 {$2$}
		 {$\cos 2x$}
	 \loigiai{
		Ta có $\dfrac{\dfrac{\sin \alpha}{\cos \alpha} - \dfrac{\cos \alpha}{\sin \alpha}}{\dfrac{\sin \alpha}{\cos \alpha} + \dfrac{\cos \alpha}{\sin \alpha}} = \dfrac{\dfrac{\sin^2 \alpha -  \cos^2 \alpha}{\sin \alpha\cdot\cos \alpha}}{\dfrac{\sin^2 \alpha + \cos^2 \alpha}{\sin \alpha\cdot\cos \alpha}} = \dfrac{\sin^2 \alpha - \cos^2 \alpha}{\sin^2 \alpha + \cos^2 \alpha} = \sin^2 \alpha - \cos^2 \alpha = -\cos 2\alpha$. \\
		Do đó $A = -\cos 2\alpha +  \cos 2\alpha = 0$.
	 }
\end{ex}

% nguồn: Ok_Toan11_TracNghiem_Lan2
\begin{ex}%[1D1H3-3]%[ma:f404ca0f]
	 Cho $\cos \alpha =-\dfrac{3}{5}$; $\dfrac{\pi}{2}<\alpha<\pi$ thì $\sin 2\alpha$ bằng
	\choice 
		{\True $-\dfrac {24}{25}$} 
		{$\dfrac {24}{25}$} 
		{$\dfrac  {4}{5}$}
		{$-\dfrac  {4}{5}$}
	\loigiai{Vì $\dfrac {\pi}{2}<\alpha<\pi$ nên $\sin \alpha>0$; $\cos  \alpha =\dfrac {-3}{5}$.\\
		Ta có $\sin^2\alpha+\cos^2\alpha=1\Rightarrow \sin  \alpha =\dfrac {4}{5}$.\\
		$\sin 2\alpha =2\sin \alpha \cos \alpha =2\cdot \dfrac {4}{5}\cdot \dfrac {-3}{5}=\dfrac {-24}{25}$. 
	 }
\end{ex}

% nguồn: Ok_Toan11_TracNghiem_Lan2
\begin{ex}%[1D1H3-3]%[ma:93670982]
	Biểu thức $\dfrac{1+\sin 4\alpha-\cos 4\alpha}{1+\sin 4\alpha+\cos 4\alpha}$ có kết quả rút gọn  bằng
	 \choice
		{$\sin  2\alpha$}
		{$\cos 2\alpha$} 
		{\True $\tan 2\alpha$} 
		{$\cot 2\alpha$} 
	\loigiai{ 
		\begin{eqnarray*}
			 && \dfrac{1+\sin 4\alpha-\cos 4\alpha}{1+\sin 4\alpha+\cos 4\alpha}\\
			&=&\dfrac{2\sin ^22\alpha+2\sin 2\alpha  \cos 2\alpha}{2\cos ^22\alpha+2\sin 2\alpha \cos 2\alpha}\\
			&=&\dfrac{2\sin 2\alpha \left(\sin 2\alpha+\cos 2\alpha  \right)}{2\cos 2\alpha \left(\sin 2\alpha+\cos 2\alpha \right)}\\
			 &=&\tan 2\alpha.
		\end{eqnarray*}
	 }
\end{ex}

% nguồn: Ok_Toan11_TracNghiem_Lan2
\begin{ex}%[1D1H3-3]%[ma:fea360ed]
	Rút gọn biểu thức $A=\sin \alpha\cdot \cos ^5\alpha-\sin ^5\alpha\cdot  \cos \alpha$.
	 \choice
		{$\dfrac{1}{2} \sin  2\alpha$}
		 {$-\dfrac{1}{2} \sin 4\alpha$}
		{$\dfrac{3}{4} \sin 4\alpha$} 
		{\True  $\dfrac{1}{4} \sin 4\alpha$}
	 \loigiai{
		 Ta có
		\begin{eqnarray*}
			\sin \alpha\cdot \cos ^5\alpha-\sin ^5\alpha\cdot \cos \alpha &= &\sin  \alpha\cdot \cos \alpha \left(\cos ^4\alpha-\sin ^4\alpha \right) \\
			 &=&\dfrac{1}{2} \sin 2\alpha \left(\cos ^2\alpha-\sin ^2\alpha \right)\left(\cos ^2\alpha+\sin ^2\alpha \right)\\
			&=&\dfrac{1}{2} \sin 2\alpha  \left(\cos ^2\alpha-\sin ^2\alpha \right)\\
			&= &\dfrac{1}{2} \sin 2\alpha \cos  2\alpha=\dfrac{1}{4} \sin 4\alpha.
		\end{eqnarray*}
	 }
\end{ex}

% nguồn: Ok_Toan11_TracNghiem_Lan2
\begin{ex}%[1D1H3-3]%[ma:2970d9dd]
	 Rút gọn biểu thức $A=\dfrac{1-\sin a-\cos 2a}{\sin 2a-\cos a}$.
	\choice 
		{$1$} 
		{\True  $\tan \alpha$}
		 {$\dfrac{5}{2}$}
		{$2\tan \alpha$} 
	\loigiai{ 
		Ta có $A=\dfrac{1-\sin a+2\sin ^2a-1}{2\sin a.\cos a-\cos a}=\dfrac{\sin a\left(2\sin a-1\right)}{\cos a\left(2\sin a-1\right)}=\dfrac{\sin a}{\cos a}=\tan a$. 
	 }
\end{ex}

% nguồn: Ok_Toan11_TracNghiem_Lan2
\begin{ex}%[1D1H3-3]%[ma:fd690f9d]
	Khi  $\alpha=\dfrac{\pi}{6}$ thì biểu thức $\dfrac{\sin ^22\alpha+4\sin ^4\alpha-4\sin ^2\alpha.\cos ^2\alpha}{4-\sin ^22\alpha-4\sin ^2\alpha}$ có giá trị bằng
	\choice 
		 {$\dfrac{1}{3}$}
		 {$\dfrac{1}{6}$}
		{\True  $\dfrac{1}{9}$}
		 {$\dfrac{1}{12}$}
	\loigiai{ 
		\begin{eqnarray*}
			&& \dfrac{\sin ^22\alpha+4\sin ^4\alpha-4\sin ^2\alpha\cdot \cos ^2\alpha}{4-\sin  ^22\alpha-4\sin ^2\alpha}\\
			&=&\dfrac{4\sin  ^4\alpha}{4(1-\sin ^2\alpha)-4\sin ^2\alpha\cdot \cos ^2\alpha} \\
			 &=&\dfrac{\sin ^4\alpha}{\cos ^2\alpha (1-\sin ^2\alpha)}
			\\&=&\dfrac{\sin ^4\alpha}{\cos  ^4\alpha}\\
			&=&\tan ^4a. 
		\end{eqnarray*}
	Khi $\alpha = \dfrac{\pi }{6}$ thì biểu  thức có giá trị bằng $\tan ^4 \dfrac{\pi}{6}=\dfrac{1}{9}$. }
\end{ex}

% nguồn: Ok_Toan11_TracNghiem_Lan2
\begin{ex}%[1D1H3-3]%[ma:e0da61d0]
	Biết $\sin x=\dfrac{1}{3}$ và $90^{\circ}<  x < 180^{\circ}$ thì biểu thức $\dfrac{1+\sin 2x+\cos 2x}{1+\sin 2x-\cos 2x}$ có giá trị bằng
	 \choice
		 {$2\sqrt{2}$}
		 {$\dfrac{1}{2\sqrt{2}}$}
		{\True  $-2\sqrt{2}$}
		 {$\dfrac{-1}{2\sqrt{2}}$}
	 \loigiai{
		Ta có $\sin x=\dfrac{1}{3}$  và $90^{\circ}< x < 180^{\circ}$ nên $\cos x <0$.\\
		Khi  đó $\cos x =-\sqrt{1-\sin^2 x}=\dfrac{-2\sqrt{2}}{3}.$\\
		Suy ra $ \sin  2x=2\cdot \sin x\cdot \cos x=\dfrac{-4\sqrt{2}}{9}$ và $ \cos 2x=1-2\sin ^2x=\dfrac{7}{9}.$\\
		Thay vào biểu  thức ta được $\dfrac{1+\sin 2x+\cos 2x}{1+\sin 2x-\cos 2x}$ $=\dfrac{1-\dfrac{4\sqrt{2}}{9}+\dfrac{7}{9}}{1-\dfrac{4\sqrt{2}}{9}-\dfrac{7}{9}}=-2\sqrt{2}.$
	 }
\end{ex}

% nguồn: Ok_Toan11_TracNghiem_Lan2
\begin{ex}%[1D1H3-3]%[ma:6ce93a6f]
	Tính giá trị của $M=\cos  ^415^\circ-\sin ^415^\circ $.
	\choice 
		{$M=1$} 
		{\True  $M=\dfrac{\sqrt{3}}{2}$}
		{$M=\dfrac{1}{4}$} 
		{$M=0$} 
	\loigiai{ 
		\begin{eqnarray*}
			 M&=&\cos ^415^\circ-\sin ^415^\circ\\
			&=&\left(\cos ^215^\circ \right)^2-\left(\sin  ^215^\circ \right)^2\\
			&=&\left(\cos ^215^\circ-\sin ^215^\circ \right)\left(\cos ^215^\circ+\sin ^215^\circ \right) 
			\\ &=&\cos ^215^\circ-\sin ^215^\circ 
			 \\ &=&\cos \left(2\cdot 15^\circ \right)\\
			&=&  \cos 30^\circ=\dfrac{\sqrt{3}}{2}.
		\end{eqnarray*}
	} 
\end{ex}

% nguồn: Ok_TracNghiemLop11_Lan1
\begin{ex}%[1D1H3-3]%[ma:1e6bdce8]
	Cho $\cos \alpha=-\dfrac{2}{3}$. Khi đó, $\cos  2\alpha$ bằng
	 \choice
		{$-\dfrac{2}{3}$} 
		 {$\dfrac{4}{3}$}
		 {$-\dfrac{4}{3}$}
		 {\True $-\dfrac{1}{9}$}
	 \loigiai{
		Ta  có $\cos 2\alpha=2\cos^2 \alpha-1=-\dfrac{1}{9}$.
	} 
\end{ex}


%%%%%%%%%%%%%%%% ĐÚNG/SAI %%%%%%%%%%%%%%%%

% nguồn: Ok_dungsai_lop11_Lan2
\begin{ex}%[1D1N3-3]%[ma:bba23c1f]
	Cho biết $\sin \alpha = \dfrac{1}{3}$  và $\dfrac{\pi}{2} < \alpha < \pi$.
	\choiceTF 
		{\True $\cos \alpha  = -\dfrac{2\sqrt{2}}{3}$}
		{\True $\sin 2\alpha  = -\dfrac{4\sqrt{2}}{9}$}
		 {\True $\cos 2\alpha = \dfrac{7}{9}$}
		{$\cot 2\alpha = \dfrac{7\sqrt{2}}{8}$} 
	\loigiai{ 
		\begin{itemchoice}
			\itemch Ta có $\dfrac{\pi}{2}  < \alpha < \pi$ nên $\cos\alpha < 0$.\\
			Ta có $\cos^2 \alpha = 1 - \sin^2 \alpha =  1 - \dfrac{1}{9} = \dfrac{8}{9} \Rightarrow \cos\alpha = -\dfrac{2\sqrt{2}}{3}$.
			\itemch  Ta có $\sin 2\alpha = 2\sin\alpha\cos\alpha = 2 \cdot \dfrac{1}{3} \cdot \left(-\dfrac{2\sqrt{2}}{3}\right) = -\dfrac{4\sqrt{2}}{9}$.
			\itemch Ta có $\cos 2\alpha = 1 - 2\sin^2 \alpha = 1 - 2\left(\dfrac{1}{3}\right)^2 = 1 - 2 \cdot  \dfrac{1}{9} = \dfrac{7}{9}$.
			\itemch Ta có $\tan 2\alpha = \dfrac{\sin 2\alpha}{\cos 2\alpha} =  \frac{-\dfrac{4\sqrt{2}}{9}}{\dfrac{7}{9}} = -\dfrac{4\sqrt{2}}{7} \Rightarrow \cot 2\alpha = \dfrac{1}{\tan 2\alpha} = -\dfrac{7}{4\sqrt{2}} = -\dfrac{7\sqrt{2}}{8}$.
		\end{itemchoice}
	} 
\end{ex}

% nguồn: Ok_dungsai_lop11_Lan1
\begin{ex}%[1D1H3-3]%[ma:5cfb5634]
 Mỗi kết  quả dưới đây đúng hay sai?
  \choiceTF
   {\True $\cos^4x-\sin^4x=\cos2x$}
  {Hàm số $f(x)=\sin x$ là hàm số  chẵn}
   {Phương trình $\cos 2x=0 $ có một nghiệm là $x=\pi$}
  {\True Phương trình $\tan 2x=1 $ có  một nghiệm là $x=\dfrac{\pi}{8}$}
  \loigiai{
  \begin{itemchoice}
    \itemch Đúng. $\cos^4x-\sin^4x=(\cos^2x+\sin^2x)(\cos^2x-\sin^2x)=\cos^2x-\sin^2x=\cos2x$.
   \itemch  Hàm số $f(x)=\sin x$ là hàm số lẻ.
    \itemch
    Do $\cos 2\pi=1$ nên $x=\pi$ không là nghiệm của phương trình $\cos 2x=0 $.
    \itemch
   Do $\tan 2\dfrac{\pi}{8}=1$  nên $x=\pi$ là nghiệm của phương trình $\tan 2x=0 $.
  \end{itemchoice}
  }
\end{ex}

% nguồn: Ok_dungsai_lop11_Lan1
\begin{ex}%[1D1H3-3]%[ma:13663e2b]
 Cho $\sin\alpha=\dfrac{3}{5}$ và $\cos \alpha=-\dfrac{4}{5}$. Mỗi  khẳng định dưới đây đúng hay sai?
  \choiceTF
   {$\sin2\alpha=\dfrac{6}{5}$}
   {\True $\cos2\alpha=\dfrac{7}{25}$}
   {\True $\tan2\alpha=-\dfrac{24}{7}$}
  {\True  $\cot2\alpha=-\dfrac{7}{24}$}
 \loigiai{ 
  \begin{itemchoice}
   \itemch $\sin2\alpha=2\sin\alpha\cos\alpha=2\cdot\dfrac{3}{5}\cdot\left(-\dfrac{4}{5}\right)=-\dfrac{24}{25}$. 
   \itemch  $\cos2\alpha=1-2\sin^2\alpha=1-2\left(\dfrac{3}{5}\right)^2=\dfrac{7}{25}$.
   \itemch  $\tan2\alpha=\dfrac{\sin2\alpha}{\cos2\alpha}=-\dfrac{24}{25}:\dfrac{7}{25}=-\dfrac{24}{7}.$
    \itemch $\cot2\alpha=\dfrac{1}{\tan2\alpha}=-\dfrac{7}{24}.$
  \end{itemchoice}
  }
\end{ex}

% nguồn: Ok_dungsai_lop11_Lan1
\begin{ex}%[1D1H3-3]%[ma:a07a5d9d]
 Cho $\sin\alpha=\dfrac{3}{5}$ và $\cos \alpha=-\dfrac{4}{5}$.  Mỗi khẳng định dưới đây đúng hay sai?
  \choiceTF
  {\True  $\cos2\alpha=\dfrac{7}{25}$}
  {$\sin2\alpha=\dfrac{6}{5}$} 
  {\True  $\tan2\alpha=-\dfrac{24}{7}$}
  {\True  $\cot2\alpha=-\dfrac{7}{24}$}
 \loigiai{ 
  \begin{itemchoice}
    \itemch $\cos2\alpha=1-2\sin^2\alpha=1-2\left(\dfrac{3}{5}\right)^2=\dfrac{7}{25}$.
    \itemch $\sin2\alpha=2\sin\alpha\cos\alpha=2\cdot\dfrac{3}{5}\cdot\left(-\dfrac{4}{5}\right)=-\dfrac{24}{25}$.
   \itemch $\tan2\alpha=\dfrac{\sin2\alpha}{\cos2\alpha}=-\dfrac{24}{25}:\dfrac{7}{25}=-\dfrac{24}{7}.$ 
    \itemch $\cot2\alpha=\dfrac{1}{\tan2\alpha}=-\dfrac{7}{24}.$
  \end{itemchoice}
 } 
\end{ex}

% nguồn: Ok_dungsai_lop11_Lan1
\begin{ex}%[1D1H3-3]%[ma:da9ec801]
 Các mệnh đề  sau \textbf{đúng} hay \textbf{sai}?
  \choiceTF
  {\True Biết  $\sin \alpha=\dfrac{1}{3}$ và $\dfrac{\pi}{2}<\alpha<\pi$ khi đó $\cos \alpha=-\dfrac{2 \sqrt{2}}{3}$}
  {\True Giá trị của hàm số $f(x)=\tan 2 x-1$ tại $x=\dfrac{\pi}{8}$  bằng $0$}
  {Phương trình $\sin2x=-\dfrac{1}{2}$ tương đương $\sin2x=\sin\dfrac{\pi}{6}$} 
  {\True Phương trình $\sin2x=-2$ vô  nghiệm}
 \loigiai{ 
  \begin{itemchoice}
    \itemch \textbf{Đúng}.\\
   Vì $\dfrac{\pi}{2}<\alpha<\pi$ nên $\cos  \alpha<0$. \\
    Suy ra
   $\cos ^2 \alpha=1-\sin ^2 \alpha=1-\dfrac{1}{9}=\dfrac{8}{9} \Rightarrow \cos  \alpha=-\dfrac{2 \sqrt{2}}{3}$.
    \itemch \textbf{Đúng}.\\
   Vì $f\left(\dfrac{\pi}{8}\right)=\tan \left(2 \cdot \dfrac{\pi}{8}\right)-1=1-1=0  $.
    \itemch \textbf{Sai}.\\
    Vì
   \begin{align*}
     \sin2x=-\dfrac{1}{2}&\Leftrightarrow\sin2x=\sin\left(-\dfrac{\pi}{6}\right)\\
     &\Leftrightarrow\hoac{&2x=-\dfrac{\pi}{6}+k2\pi\\&2x=\dfrac{7\pi}{6}+k2\pi}\\
     &\Leftrightarrow\hoac{&x=-\dfrac{\pi}{12}+k\pi\\&x=\dfrac{7\pi}{12}+k\pi}(k\in\mathbb{Z}).
   \end{align*}
   \itemch \textbf{Đúng}.\\ 
   Vì $-2 \notin  [-1;1]$ nên phương trình $\sin2x=-2$ vô nghiệm.
  \end{itemchoice}
  }
\end{ex}

% nguồn: Ok_dungsai_lop11_Lan1
\begin{ex}%[1D1H3-3]%[ma:6a38a1d1]
 Mỗi kết quả  dưới đây đúng hay sai?
  \choiceTF
  {$\dfrac{\cos2x}{\sin  x+\cos x}=\sin x-\cos x$}
   {\True $\dfrac{\cos^2x-\sin^2x}{\cos x-\sin x}=\cos x+\sin x$}
  {\True  $\cos^4x-\sin^4x=\cos2x$}
   {$\dfrac{\cos2x+1}{1-\sin^2x}=1$}
 \loigiai{ 
  \begin{itemchoice}
   \itemch $\dfrac{\cos2x}{\sin  x+\cos x}=\dfrac{\cos^2x-\sin^2x}{\sin x+\cos x}=\dfrac{(\cos x+\sin x)(\cos x-\sin x)}{\sin x+\cos x}=\cos x-\sin x$.
   \itemch $\dfrac{\cos^2x-\sin^2x}{\cos x-\sin x}=\dfrac{(\cos x+\sin x)(\cos x-\sin x)}{\cos  x-\sin x}=\cos x+\sin x$.
   \itemch $\cos^4x-\sin^4x=(\cos^2x+\sin^2x)(\cos^2x-\sin^2x)=\cos^2x-\sin^2x=\cos2x$. 
   \itemch $\dfrac{\cos2x+1}{1-\sin^2x}=\dfrac{(2\cos^2x-1)+1}{\cos^2x}=\dfrac{2\cos^2x}{\cos^2x}=2$. 
  \end{itemchoice}
  }
\end{ex}

% nguồn: Ok_dungsai_lop11_Lan2
\begin{ex}%[1D1H3-3]%[ma:46e48e49]
	Cho $\cos  x=\dfrac{1}{5}$, $\dfrac{\pi}{2}< x <\pi $. Khi đó
	\choiceTF 
		 {$\sin\dfrac{x}{2}=\dfrac{\sqrt{10}}{4}$}
		 {$\cos\dfrac{x}{2}=\dfrac{\sqrt{15}}{4}$}
		{\True  $\tan\dfrac{x}{2}=\dfrac{\sqrt 6}{3}$}
		{\True $\cot\dfrac{x}{2}=\dfrac{\sqrt 6}{2}$} 
	\loigiai{ 
		\begin{itemchoice}
			\itemch \textbf{Sai.}  \\
			$\dfrac{\pi}{2}<  x <\pi\Rightarrow\dfrac{\pi}{4}<\dfrac{x}{2}<\dfrac{\pi}{2}\Rightarrow\sin\dfrac{x}{2}> 0.\\
			\sin\dfrac{x}{2}=\sqrt{\dfrac{1-\cos x}{2}}=\sqrt{\dfrac{1-\dfrac{1}{5}}{2}}=\dfrac{\sqrt{10}}{5}$. 
			 \itemch \textbf{Sai.} \\
			 $\dfrac{\pi}{2}< x <\pi\Rightarrow\dfrac{\pi}{4}<\dfrac{x}{2}<\dfrac{\pi}{2}\Rightarrow\cos\dfrac{x}{2}> 0.\\
			\cos\dfrac{x}{2}=\sqrt{\dfrac{1+\cos  x}{2}}=\sqrt{\dfrac{1+\dfrac{1}{5}}{2}}=\dfrac{\sqrt{15}}{5}$.
			 \itemch \textbf{Đúng.} \\
			 $\tan\dfrac{x}{2}=\dfrac{\sin\dfrac{x}{2}}{\cos\dfrac{x}{2}}=\dfrac{\dfrac{\sqrt{10}}{5}}{\dfrac{\sqrt{15}}{5}}=\dfrac{\sqrt 6}{3}$.
			 \itemch \textbf{Đúng.} \\
			 $\cot\dfrac{x}{2}=\dfrac{\cos\dfrac{x}{2}}{\sin\dfrac{x}{2}}=\dfrac{\dfrac{\sqrt{15}}{5}}{\dfrac{\sqrt{10}}{5}}=\dfrac{\sqrt 6}{2}$.
		\end{itemchoice}
	} 
\end{ex}

% nguồn: Ok_dungsai_lop11_Lan2
\begin{ex}%[1D1H3-3]%[ma:2c2af5b2]
	 Cho biết $\cos 2\alpha = -\dfrac{1}{4}$ và $\pi < \alpha < \dfrac{3\pi}{2}$.
	 \choiceTF
		{\True $\sin \alpha  < 0$, $\cos \alpha < 0$}
		{$\sin \alpha  = \dfrac{\sqrt{10}}{4}$}
		{$\cos \alpha =  \dfrac{\sqrt{6}}{4}$}
		{\True $\cot \alpha =  \dfrac{\sqrt{15}}{5}$}
	 \loigiai{
		\begin{itemchoice}
			\itemch Vì $\pi < \alpha <  \dfrac{3\pi}{2}$ nên $\sin\alpha < 0, \cos\alpha < 0$.
			\itemch Ta có $\cos 2\alpha = -\dfrac{1}{4} \Rightarrow 1 - 2\sin^2\alpha =  -\dfrac{1}{4} \Rightarrow \sin^2\alpha = \dfrac{5}{8} \Rightarrow \sin\alpha = -\dfrac{\sqrt{10}}{4}$.
			\itemch Ta có $\cos\alpha  = -\sqrt{1-\sin^2\alpha} = -\sqrt{1-\dfrac{10}{16}} = -\dfrac{\sqrt{6}}{4}$.
			\itemch Ta có $\tan\alpha = \dfrac{\sin\alpha}{\cos\alpha} =  \dfrac{\sqrt{15}}{3} \Rightarrow \cot\alpha = \dfrac{1}{\tan\alpha} = \dfrac{\sqrt{15}}{5}$.
		\end{itemchoice}
	} 
\end{ex}

% nguồn: Ok_dungsai_lop11_Lan2
\begin{ex}%[1D1H3-3]%[ma:9787f66b]
	Cho $\cos\alpha=\dfrac{2}{5},2\pi <\alpha <\dfrac{5\pi}{2}$,  khi đó
	\choiceTF 
		{\True $\sin\alpha=\dfrac{\sqrt{21}}{5}$} 
		{\True  $\sin 2\alpha=\dfrac{4\sqrt{21}}{25}$}
		{\True $\cos 2\alpha=-\dfrac{17}{25}$} 
		 {\True $\tan 2\alpha=\dfrac{-4\sqrt{21}}{17}$}
	 \loigiai{
		\begin{itemchoice}
			\itemch \textbf{Đúng.} \\ 
			Ta có  ${\sin ^2}\alpha=1-\cos ^2\alpha=1-\left(\dfrac{2}{5}\right)^2=\dfrac{21}{25}\Rightarrow\sin\alpha=\pm\dfrac{\sqrt{21}}{5}$.\\
			Vì $2\pi <\alpha  <\dfrac{5\pi}{2}$ nên $\sin\alpha=\dfrac{\sqrt{21}}{5}$.
			\itemch  \textbf{Đúng.} \\
			$\sin 2\alpha=2\sin\alpha\cos\alpha=2\cdot\dfrac{\sqrt{21}}{5}\cdot\dfrac{2}{5}=\dfrac{4\sqrt{21}}{25}$. 
			 \itemch \textbf{Đúng.} \\
			$\cos  2\alpha=\cos ^2\alpha-\sin ^2\alpha=\left(\dfrac{2}{5}\right)^2-\left(\dfrac{\sqrt{21}}{5}\right)^2=-\dfrac{17}{25}$.
			\itemch  \textbf{Đúng.} \\
			$\tan  2\alpha=\dfrac{\sin 2\alpha}{\cos 2\alpha}=\dfrac{4\sqrt{21}}{25}:\dfrac{-17}{25}=\dfrac{-4\sqrt{21}}{17}.$
		\end{itemchoice}
	 }
\end{ex}

% nguồn: Ok_dungsai_lop11_Lan2
\begin{ex}%[1D1H3-3]%[ma:11071dcc]
	Cho $\sin x=\dfrac{1}{5},\dfrac{\pi}{2}< x <\pi $. Khi  đó
	 \choiceTF
		{$\sin  2x=\dfrac{4\sqrt 6}{5}$}
		{\True  $\cos 2x=\dfrac{23}{25}$}
		{$\tan 2x=\dfrac{20\sqrt  6}{3}$}
		{$\cot 2x=\dfrac{23\sqrt 6}{120}$} 
	\loigiai{ 
		\begin{itemchoice}
			\itemch  \textbf{Sai.} \\
			$\dfrac{\pi}{2}< x <\pi\Rightarrow\dfrac{2\pi}{2}< 2x < 2\pi\Rightarrow\pi < 2x < 2\pi\Rightarrow\sin 2x < 0$.\\ 
			 $\dfrac{\pi}{2}< x <\pi\Rightarrow\cos x < 0.\\
			\sin x=\dfrac{1}{5}\Rightarrow\cos  x=-\sqrt{1-\sin^2x}=-\dfrac{2\sqrt 6}{5}$.\\
			$\sin 2x=2\sin x\cdot\cos  x=2\cdot\dfrac{1}{5}\cdot\left(-\dfrac{2\sqrt 6}{5}\right)=-\dfrac{4\sqrt 6}{5}$.
			\itemch  \textbf{Đúng.} \\
			$\dfrac{\pi}{2}< x <\pi\Rightarrow\dfrac{2\pi}{2}< 2x <  2\pi\Rightarrow\pi < 2x < 2\pi\Rightarrow\cos 2x > 0$.\\
			$\cos  2x=1-2\sin ^2x=1-2\cdot\dfrac{1}{25}=\dfrac{23}{25}$.
			\itemch \textbf{Sai.}  \\
			$\tan 2x=\dfrac{\sin 2x}{\cos  2x}=\dfrac{-\dfrac{4\sqrt 6}{5}}{\dfrac{23}{25}}=-\dfrac{20\sqrt 6}{3}$.
			 \itemch \textbf{Sai.} \\
			$\cot 2x=\dfrac{\cos 2x}{\sin 2x}=\dfrac{\dfrac{23}{25}}{-\dfrac{4\sqrt 6}{5}}=-\dfrac{23\sqrt  6}{120}$.
		\end{itemchoice}
	} 
\end{ex}

% nguồn: Ok_dungsai_lop11_Lan2
\begin{ex}%[1D1H3-3]%[ma:f02f675a]
	Biết $\cos 2\alpha=\dfrac{5}{9}$, $0^\circ < \alpha < 90^\circ$. Khi đó 
	\choiceTF 
		{\True  $\sin\alpha=\dfrac{\sqrt{28}}{9}$}
		{\True  $\cos\alpha=\dfrac{\sqrt{53}}{9}$}
		{$\tan\alpha=\dfrac{\sqrt{371}}{53}$} 
		 {\True $\cot\alpha=\dfrac{\sqrt{371}}{14}$}
	 \loigiai{
		\begin{itemchoice}
			 \itemch \textbf{Đúng.} \\
			$\sin ^2\alpha=\dfrac{1-\cos  2\alpha}{2}=\dfrac{1-\left(\dfrac{5}{9}\right)^2}{2}=\dfrac{28}{81}\Rightarrow\sin\alpha=\pm\dfrac{\sqrt{28}}{9}$.\\
			Vì $0^\circ < \alpha <  90^\circ$ nên $\sin\alpha=\dfrac{\sqrt{28}}{9}$.
			 \itemch \textbf{Đúng.} \\
			$\cos ^2\alpha=1-\sin  ^2\alpha=1-\left(\dfrac{\sqrt{28}}{9}\right)^2=\dfrac{53}{81}\Rightarrow\cos\alpha=\pm\dfrac{\sqrt{53}}{81}$.\\
			Vì $0^\circ <  \alpha < 90^\circ$ nên $\cos\alpha=\dfrac{\sqrt{53}}{9}$.
			\itemch  \textbf{Sai.} \\
			$\tan\alpha=\dfrac{\sin\alpha}{\cos\alpha}=\dfrac{2\sqrt{371}}{53}$. 
			\itemch \textbf{Đúng.}  \\
			$\cot\alpha=\dfrac{1}{\tan\alpha}=\dfrac{\sqrt{371}}{14}$. 
		\end{itemchoice}
	} 
\end{ex}


%%%%%%%%%%%%%%%% TRẢ LỜI NGẮN %%%%%%%%%%%%%%%%

% nguồn: Ok_TraLoiNgan_Lop11_Lan2
\begin{ex}%[1D1H3-3]%[ma:e72f4cc2]
	Cho $\sin x=\dfrac{1}{5}$, $\dfrac{\pi}{2}< x <\pi $. Tính $\cot 2x$ (Kết  quả làm tròn đến hàng phần chục).
	 \shortans{-0,5}
	 \loigiai{
		\begin{itemize}
			\item $\cos  2x=1-2\sin ^2x=1-2\cdot\dfrac{1}{25}=\dfrac{23}{25}.$
			\item $\dfrac{\pi}{2}<  x <\pi\Rightarrow\cos x < 0$.\\
			 $\sin x=\dfrac{1}{5}\Rightarrow\cos x=-\sqrt{1-\sin^2x}=-\dfrac{2\sqrt 6}{5}$.
			\item $\sin 2x=2\sin x\cdot\cos  x=2\cdot\dfrac{1}{5}\cdot\left(-\dfrac{2\sqrt 6}{5}\right)=-\dfrac{4\sqrt 6}{5}$.
			 \item $\cot 2x=\dfrac{\cos 2x}{\sin 2x}=\dfrac{\dfrac{23}{25}}{-\dfrac{4\sqrt 6}{5}}=-\dfrac{23\sqrt 6}{120}\approx -0{,}5$.
		\end{itemize}
	} 
\end{ex}

% nguồn: Ok_TraLoiNgan_Lop11_Lan2
\begin{ex}%[1D1H3-3]%[ma:e38d5298]
	Cho tam giác $ABC$. Biết $\sin A+\sin  B+\sin C=k\cos\dfrac{A}{2}\cos\dfrac{B}{2}\cos\dfrac{C}{2}$, tìm $k$?
	 \shortans{4}
	 \loigiai{
		\begin{eqnarray*}
			 VT&=&\sin A+\sin B+\sin C\\
			&=&2\sin\dfrac{A+B}{2}\cos\dfrac{A-B}{2}+2\sin\dfrac{C}{2}\cos\dfrac{C}{2}\\ 
			&=&2\sin\dfrac{\pi-C}{2}\cos\dfrac{A-B}{2}+2\cos\dfrac{\pi-C}{2}\cos\dfrac{C}{2}\\ 
			 &=&2\cos\dfrac{C}{2}\cos\dfrac{A-B}{2}+2\cos\dfrac{A+B}{2}\cos\dfrac{C}{2}\\
			 &=&2\cos\dfrac{C}{2}\left(\cos\dfrac{A-B}{2}+\cos\dfrac{A+B}{2}\right)\\
			&=&2\cos\dfrac{C}{2}\cdot  2\cos\dfrac{A-B+A+B}{4}\cdot\cos\dfrac{A-B-A-B}{4}\\
			 &=&4\cos\dfrac{C}{2}\cos\dfrac{A}{2}\cos\dfrac{-B}{2}=4\cos\dfrac{A}{2}\cos\dfrac{B}{2}\cos\dfrac{C}{2}.
		\end{eqnarray*}
		Vậy $\sin A+\sin B+\sin  C=4\cos\dfrac{A}{2}\cos\dfrac{B}{2}\cos\dfrac{C}{2}$. .
	} 
\end{ex}

% nguồn: Ok_TraLoiNgan_Lop11_Lan1
\begin{ex}%[1D1V3-3]%[ma:39be03a2]
	Số giờ có ánh sáng mặt trời của một thành phố $A$ ở vĩ độ $40^{\circ}$ bắc trong ngày  thứ $t$ của một năm không nhuận được cho bởi hàm số
	 $$
	d(t)=3\sin \left[\frac{\pi}{182}(t-80)\right]+12 \text {  với } t \in \mathbb{Z} \text { và } 0<t \leq 365.
	$$ 
	Thành phố $A$ có đúng 12 giờ có ánh sáng mặt  trời vào ngày nào đầu tiên trong năm?
	 \shortans{82}
	\loigiai{ 
		Ta có  phương trình
		\begin{eqnarray*}
			&&3\sin \left[\dfrac{\pi}{182}(t-80)\right]+12=12\\ 
			 &\Leftrightarrow & \sin \left[\dfrac{\pi}{182}(t-80)\right]=0 \\
			 &\Leftrightarrow &
			\dfrac{\pi}{182}(t-80)=k\pi \Leftrightarrow  t=80+182k
		\end{eqnarray*}
		Với $k,t \in \mathbb{Z}$ và $0<t \leq 365$, ta tìm được $k=0$ và $k=1$  thỏa.\\
		Suy ra thành phố $A$ có 12  giờ ánh sáng vào ngày thứ $80$ và ngày thứ $262$ của năm.
	} 
\end{ex}

% nguồn: Ok_TraLoiNgan_Lop11_Lan2
\begin{ex}%[1D1V3-3]%[ma:f5d720ee]
	Cho $\sin \alpha+\cos \alpha=m$. Có bao nhiêu  giá trị của tham số $m$ để $\sin 2\alpha=-\dfrac{3}{4}$?\shortans{2}
	 \loigiai{
		Ta có $\sin \alpha+\cos \alpha=\sqrt{2} \left(\dfrac{\sqrt{2}}{2} \sin \alpha+\dfrac{\sqrt{2}}{2} \cos \alpha  \right)=\sqrt{2} \sin \left(\alpha+\dfrac{\pi}{4} \right)$.\\
		Vì $-1\le \sin \left(\alpha+\dfrac{\pi}{4} \right)\le 1$ nên $-\sqrt{2} \le \sin \alpha+\cos \alpha \le \sqrt{2}$.\\ 
		Suy ra $-\sqrt{2} \le  m\le \sqrt{2}$.\\
		Ta lại có $(\sin \alpha+\cos \alpha)^2=\sin  ^2\alpha+2\sin \alpha \cos \alpha+\cos ^2\alpha=1+\sin 2\alpha$.\\
		Suy ra  $\sin 2\alpha=(\sin \alpha+\cos \alpha)^2-1=m^2-1$.\\
		Khi đó, $\sin  2\alpha=-\dfrac{3}{4}$ hay $m^2-1=-\dfrac{3}{4}$.\\
		Suy ra $m=\dfrac{1}{2}$ hoặc $m=-\dfrac{1}{2}$ (thoả  mãn điều kiện).
	} 
\end{ex}


%%%%%%%%%%%%%%%% TỰ LUẬN %%%%%%%%%%%%%%%%

% nguồn: Ok_TuanLuan_Toan11_Lan2
\begin{ex}%[1D1H3-3]%[ma:a25d3b65]
	Tính  $A = \cos^4 \dfrac{\pi}{12} - \sin^4 \dfrac{\pi}{12}$.
	 \loigiai{
		Ta có $A = \left(\cos^2 \dfrac{\pi}{12} - \sin^2  \dfrac{\pi}{12}\right) \left(\cos^2 \dfrac{\pi}{12} + \sin^2 \dfrac{\pi}{12}\right) = \cos^2 \dfrac{\pi}{12} - \sin^2 \dfrac{\pi}{12} = \cos \dfrac{\pi}{6} = \dfrac{\sqrt{3}}{2}$.
	} 
\end{ex}

% nguồn: Ok_TuanLuan_Toan11_Lan2
\begin{ex}%[1D1H3-3]%[ma:b5a58ffd]
	Tính $\sin 2a$, $\cos  2a$, $\tan 2a$, biết:
	\begin{listEX}[2]
		\item $\sin a = \dfrac{1}{3}$ và  $\pi < a < \dfrac{3\pi}{2}$.
		\item $\sin a + \cos a = -\dfrac{1}{2}$  và $\dfrac{\pi}{2} < a < \dfrac{3\pi}{4}$.
	\end{listEX}
	 \loigiai{
		\begin{listEX}
			\item Vì $\dfrac{\pi}{2} < a < \pi$ nên  $\cos a < 0$.\\
			Mặt khác $\sin^2 a + \cos^2 a = 1$ suy ra $\cos a = -\sqrt{1 - \sin^2 a} = -\sqrt{1  - \left(\dfrac{1}{3}\right)^2} = -\dfrac{2\sqrt{2}}{3}$.\\
			$\sin 2a  = 2\sin a \cos a = 2 \cdot \left(\dfrac{1}{3}\right) \cdot \left(-\dfrac{2\sqrt{2}}{3}\right) = -\dfrac{4\sqrt{2}}{9}$;\\ $\cos 2a = 1 - 2\sin^2 a = 1 - 2 \cdot \left(\dfrac{1}{3}\right)^2 = \dfrac{7}{9}$.\\
			Suy ra $\tan 2a =  \dfrac{\sin 2a}{\cos 2a} = \dfrac{-\dfrac{4\sqrt{2}}{9}}{\dfrac{7}{9}} = -\dfrac{4\sqrt{2}}{7}$.
			\item Ta có $(\sin a + \cos a)^2 = \left(-\dfrac{1}{2}\right)^2 \Leftrightarrow \sin^2 a + \cos^2 a + 2\sin a \cos a = \dfrac{1}{4}\Leftrightarrow 1 +  \sin 2a = \dfrac{1}{4}$.\\
			Suy ra $\sin  2a = -\dfrac{3}{4}$.\\
			Vì $\dfrac{\pi}{2} < a  < \dfrac{3\pi}{4}$ nên $\pi < 2a < \dfrac{3\pi}{2}$. Do đó $\cos 2a < 0$.\\
			Mặt khác $\sin^2 (2a) +  \cos^2 (2a) = 1$ suy ra $\cos 2a = -\sqrt{1 - \sin^2 (2a)} = -\sqrt{1 - \left(-\dfrac{3}{4}\right)^2} = -\dfrac{\sqrt{7}}{4}$.\\
			Do đó $\tan 2a = \dfrac{\sin 2a}{\cos 2a} = \dfrac{-\dfrac{3}{4}}{-\dfrac{\sqrt{7}}{4}} = \dfrac{3}{\sqrt{7}} =  \dfrac{3\sqrt{7}}{7}$.
		\end{listEX}
	 }
\end{ex}

% nguồn: Ok_TuanLuan_Toan11_Lan2
\begin{ex}%[1D1H3-3]%[ma:fef83e47]
	Cho $\cos  2x = -\dfrac{4}{5}$, với $\dfrac{\pi}{4} < x < \dfrac{\pi}{2}$. Tính $\sin x$, $\cos x$, $\sin \left(x+\dfrac{\pi}{3}\right)$, $\cos \left(2x-\dfrac{\pi}{4}\right)$.
	\loigiai{ 
		Vì $\dfrac{\pi}{4} < x < \dfrac{\pi}{2}$ nên $\sin x > 0$, $\cos x  > 0$.\\
		Áp dụng công  thức hạ bậc, ta có\\
		$\sin^2 x = \dfrac{1 - \cos 2x}{2} =  \dfrac{9}{10} \Rightarrow \sin x = \dfrac{3}{\sqrt{10}}$;
		$\cos^2 x = \dfrac{1 + \cos 2x}{2}= \dfrac{1}{10} \Rightarrow \cos x = \sqrt{\dfrac{1}{10}}  = \dfrac{1}{\sqrt{10}}$.\\
		 Theo công thức cộng, ta có\\
		$\sin \left(x+\dfrac{\pi}{3}\right) = \sin  x \cos \dfrac{\pi}{3} + \cos x \sin \dfrac{\pi}{3} = \dfrac{3}{\sqrt{10}} \cdot \dfrac{1}{2} + \dfrac{1}{\sqrt{10}} \cdot \dfrac{\sqrt{3}}{2} = \dfrac{3 + \sqrt{3}}{2\sqrt{10}} = \dfrac{3\sqrt{10} + \sqrt{30}}{20}$.\\
		$\cos \left(2x-\dfrac{\pi}{4}\right) = \cos 2x  \cos \dfrac{\pi}{4} + \sin 2x \sin \dfrac{\pi}{4}= \left(-\dfrac{4}{5}\right) \cdot \dfrac{\sqrt{2}}{2} +\dfrac{\sqrt{2}}{2}\cdot 2\cdot\dfrac{3}{\sqrt{10}} \cdot \dfrac{1}{\sqrt{10}} = -\dfrac{\sqrt{2}}{10}$.
	 }
\end{ex}

% nguồn: Ok_TuanLuan_Toan11_Lan2
\begin{ex}%[1D1V3-3]%[ma:0c130006]
	Tính giá trị  của các biểu thức sau:
	\begin{listEX}[2]
		\item $A = \dfrac{\sin \dfrac{\pi}{15} \cos \dfrac{\pi}{10} + \sin \dfrac{\pi}{10} \cos  \dfrac{\pi}{15}}{\cos \dfrac{2\pi}{15} \cos \dfrac{\pi}{5} - \sin \dfrac{2\pi}{15} \sin \dfrac{\pi}{5}}$.
		\item $B =  \sin \dfrac{\pi}{32} \cos \dfrac{\pi}{32} \cos \dfrac{\pi}{16} \cos \dfrac{\pi}{8}$.
	\end{listEX}
	 \loigiai{
		\begin{listEX}
			\item Ta có 
			$A =  \dfrac{\sin \dfrac{\pi}{15} \cos \dfrac{\pi}{10} + \sin \dfrac{\pi}{10} \cos \dfrac{\pi}{15}}{\cos \dfrac{2\pi}{15} \cos \dfrac{\pi}{5} - \sin \dfrac{2\pi}{15} \sin \dfrac{\pi}{5}}= \dfrac{\sin \left(\dfrac{\pi}{15} + \dfrac{\pi}{10}\right)}{\cos \left(\dfrac{2\pi}{15} + \dfrac{\pi}{5}\right)}= \dfrac{\sin \dfrac{\pi}{6}}{\cos \dfrac{\pi}{3}}$
			 $= \dfrac{\dfrac{1}{2}}{\dfrac{1}{2}} = 1$.
			\item Ta  có
			$B = \sin \dfrac{\pi}{32} \cos  \dfrac{\pi}{32} \cos \dfrac{\pi}{16} \cos \dfrac{\pi}{8}$\\
			$= \dfrac{1}{2} \left(2\sin \dfrac{\pi}{32} \cos \dfrac{\pi}{32}\right)  \cos \dfrac{\pi}{16} \cos \dfrac{\pi}{8}$\\
			$= \dfrac{1}{2} \sin \left(2 \cdot \dfrac{\pi}{32}\right) \cos \dfrac{\pi}{16} \cos  \dfrac{\pi}{8}$
			 $= \dfrac{1}{2} \sin \dfrac{\pi}{16} \cos \dfrac{\pi}{16} \cos \dfrac{\pi}{8}$\\
			$= \dfrac{1}{2} \cdot \dfrac{1}{2} \left(2\sin \dfrac{\pi}{16} \cos \dfrac{\pi}{16}\right) \cos  \dfrac{\pi}{8}$
			$=  \dfrac{1}{4} \sin \left(2 \cdot \dfrac{\pi}{16}\right) \cos \dfrac{\pi}{8}$\\
			$= \dfrac{1}{4}  \sin \dfrac{\pi}{8} \cos \dfrac{\pi}{8}$
			$= \dfrac{1}{4} \cdot \dfrac{1}{2} \left(2\sin \dfrac{\pi}{8} \cos  \dfrac{\pi}{8}\right)$\\
			$= \dfrac{1}{8} \sin \left(2 \cdot \dfrac{\pi}{8}\right)$ 
			$= \dfrac{1}{8}  \sin \dfrac{\pi}{4}$
			$= \dfrac{1}{8} \cdot \dfrac{\sqrt{2}}{2}  = \dfrac{\sqrt{2}}{16}$.
		\end{listEX}
	} 
\end{ex}
